% ==================================================================
% id:       lim-010
% topic:    גבולות ורציפות של פונקציות
% subtopic: חישוב גבולות של פונקציות
% solution: full
% ==================================================================
\textbf{
בכל אחד מהסעיפים הבאים,
חשבו את הגבול: }
\begin{enumerate}\bMs{2}
\item $\ds{\lim_{x\to\infty}\frac{2x^2-3x+6}{x^2-51x+1}}$
\ifsol
\begin{solbox}
מתקיים
\begin{align*}
    \frac{2x^2-3x+6}{x^2-51x+1}=\ \frac{x^2(2-\tfrac{3}{x}+\tfrac{6}{x^2})}{x^2(1-\tfrac{51}{x}+\tfrac{1}{x^2})}
    =\ \frac{2-\tfrac{3}{x}+\tfrac{6}{x^2}}{1-\tfrac{51}{x}+\tfrac{1}{x^2}},
\end{align*}
לכן
מאריתמטיקה של גבולות נקבל
\begin{align*}
    \ds{\lim_{x\to\infty}\frac{2x^2-3x+6}{x^2-51x+1}}
    =\ds{\lim_{x\to\infty}\frac{2-\tfrac{3}{x}+\tfrac{6}{x^2}}{1-\tfrac{51}{x}+\tfrac{1}{x^2}}}
    = \frac{2-0+0}{1-0+0}
    = 2,
\end{align*}
לפי אריתמטיקה של גבולות.
\end{solbox}
\fi

\item $\ds{ \lim_{x \to \infty} \lr{\frac{2x^3 - 4x^2 + 5}{x^2 + 2x +1}-2x}}$
\ifsol
\hint{אחדו את שני האיברים לשבר אחד לפני השאיפה לאינסוף.}
\begin{solbox}
מתקיים
\begin{align*}
    ,\lr{\frac{2x^3 - 4x^2 + 5}{x^2 + 2x +1}-2x}= \frac{ - 8x^2 -2x+ 5}{x^2 + 2x +1}
    = \frac{-8-\tfrac{2}{x}+\tfrac{5}{x^2}}{1+\frac{2}{x}+\frac{1}{x^2}}
\end{align*}
לכן
\begin{align*}
    .\ds{ \lim_{x \to \infty} \lr{\frac{2x^3 - 4x^2 + 5}{x^2 + 2x +1}-2x}}
    = \ds{ \lim_{x \to \infty}\frac{-8-\tfrac{2}{x}+\tfrac{5}{x^2}}{1+\frac{2}{x}+\frac{1}{x^2}}}
    = \frac{-8-0+0}{1+0+0}
    = -8
\end{align*}
הערה: בחישוב הגבול השתמשנו באריתמטיקה של גבולות: הגבול של סכום, של מכפלה ושל מנה של פונקציות מתכנסות שווה לסכום, למכפלה ולמנה של הגבולות, כל עוד הגבול של המכנה שונה מ-$0$.
\end{solbox}
\fi

\item $\ds{ \lim_{x \to \infty}\frac{x^3+3x-1}{\sqrt{9x^6+2}}}$
\ifsol
\begin{solbox}
מתקיים
\begin{align*}
    ,\frac{x^3+3x-1}{\sqrt{9x^6+2}}=\frac{x^3+3x-1}{\sqrt{x^6(9+\tfrac{2}{x^6})}}
    = \frac{x^3(1+\tfrac{3}{x^2}-\tfrac{1}{x^3})}{x^3\sqrt{9+\tfrac{2}{x^6}}}
    = \frac{1+\tfrac{3}{x^2}-\tfrac{1}{x^3}}{\sqrt{9+\tfrac{2}{x^6}}}
\end{align*}
לכן
\begin{align*}
    \ds{ \lim_{x \to \infty}\frac{x^3+3x-1}{\sqrt{9x^6+2}}}= \ds{ \lim_{x \to \infty}\frac{1+\tfrac{3}{x^2}-\tfrac{1}{x^3}}{\sqrt{9+\tfrac{2}{x^6}}}}
    = \frac{1}{3},
\end{align*}
הערה: בחישוב הגבול השתמשנו באריתמטיקה של גבולות: הגבול של סכום, של מכפלה ושל מנה של פונקציות מתכנסות שווה לסכום, למכפלה ולמנה של הגבולות, כל עוד הגבול של המכנה שונה מ-$0$ בנוסף, השורש הוא פונקציה רציפה, ולכן אפשר להעביר את הגבול אל תוך השורש.
\end{solbox}
\fi

\item $\ds{\lim_{x\to\infty}\frac{\cos(x)}{x}}$
\ifsol
\hint{מה ידוע על מכפלה של פונקציה חסומה בפונקציה ששואפת לאפס?}
\begin{solbox}
נשים לב כי
$|\cos(x)|\leq 1$
לכל
$x\in\mathbb{R}$,
כלומר,
$\cos(x)$
היא פונקציה חסומה. בנוסף
$\displaystyle\lim_{x\to\infty}\frac{1}{x}=0$.
נסיק כי
$\ds{\lim_{x\to\infty}\frac{\cos(x)}{x}}=0$
כמכפלה של פונקציה חסומה באפסה.
\end{solbox}
\fi

\item $\ds{\lim_{x\to\infty}\sin(x)\cdot\sin\left(\tfrac{1}{x}\right)}$
\ifsol
\hint{מה ידוע על מכפלה של פונקציה חסומה בפונקציה ששואפת לאפס?}
\begin{solbox}
נשים לב כי
$|\sin(x)|\leq 1$
לכל
$x\in\mathbb{R}$,
כלומר,
$\sin(x)$
היא פונקציה חסומה. בנוסף
$\displaystyle\lim_{x\to\infty}\sin\left(\frac{1}{x}\right)=0$
נסיק כי
$\ds{\lim_{x\to\infty}\sin(x)\cdot\sin\left(\tfrac{1}{x}\right)}=0$.
כמכפלה של פונקציה חסומה באפסה.
\end{solbox}
\fi

\item $\ds{ \lim_{x \to -\infty} \frac{x}{\sqrt{2x^2 + 1}}}$
\ifsol
\hint{הוציאו את $x^2$ מתוך השורש ושימו לב ש-$\sqrt{x^2}=|x|$. מה הסימן של $x$ כאן?}
\begin{solbox}
מתקיים
\begin{align*}
    ,\frac{x}{\sqrt{2x^2 + 1}}=
    \frac{x}{|x|\sqrt{2 + \frac{1}{x^2}}}
    =
    \frac{x}{-x\sqrt{2 + \frac{1}{x^2}}}
    \overset{(\ast)}{=}
    -\frac{1}{\sqrt{2 + \frac{1}{x^2}}},
\end{align*}
לכן
\begin{align*}
    .\lim_{x \to -\infty} \frac{x}{\sqrt{2x^2 + 1}}
    =
    -\lim_{x \to -\infty} \frac{1}{\sqrt{2 + \frac{1}{x^2}}}
    =\ -\frac{1}{\sqrt{2 + 0}}
    = -\frac{1}{\sqrt{2}}
\end{align*}
$(\ast)$
    יש כאן שאיפה למינוס אינסוף, לכן
    $|x|=-x$.
הערה: בחישוב הגבול השתמשנו באריתמטיקה של גבולות: הגבול של סכום, של מכפלה ושל מנה של פונקציות מתכנסות שווה לסכום, למכפלה ולמנה של הגבולות, כל עוד הגבול של המכנה שונה מ-$0$ בנוסף, השורש הוא פונקציה רציפה, ולכן אפשר להעביר את הגבול אל תוך השורש.
\end{solbox}
\fi

\item $\ds{ \lim_{x \to \infty}(\sqrt{x^2+1}-\sqrt{x^2-1})}$
\ifsol
\hint{זהו ביטוי מהצורה $\infty-\infty$ עם שורשים: הכפילו וחלקו בצמוד.}
\begin{solbox}
מתקיים
\begin{align*}
    (\sqrt{x^2+1}-\sqrt{x^2-1})&=\ (\sqrt{x^2+1}-\sqrt{x^2-1})\cdot{\color{red}\frac{\sqrt{x^2+1}+\sqrt{x^2-1}}{\sqrt{x^2+1}+\sqrt{x^2-1}}}\\
    &=\ \frac{2}{\sqrt{x^2+1}+\sqrt{x^2-1}}
\end{align*}
לכן
\begin{align*}
    \lim_{x \to \infty}(\sqrt{x^2+1}-\sqrt{x^2-1})
    &=\ \lim_{x \to \infty}\frac{2}{\sqrt{x^2+1}+\sqrt{x^2-1}}\\
    &=\ 0.
\end{align*}
\end{solbox}
\fi
\eMs
\end{enumerate}
